\documentclass[10pt]{article}
\usepackage[margin=1in]{geometry}
\usepackage{amsmath}
\usepackage{booktabs}
\usepackage{hyperref}

\title{PLA: A Transparent, Learnable Rule-Confidence Engine\\
with Measured Trace Fidelity}
\author{Seyed Masoud Hashemi Ahmadi\\
\small École de technologie supérieure (ÉTS), Montréal\\
\small \texttt{contact@AiCentralLab.com}}
\date{Working draft --- workshop/tool-track cut}

\begin{document}
\maketitle

\begin{abstract}
We present PLA, a small, fully transparent rule-confidence engine for
settings where inherent interpretability is a requirement. PLA descends
openly from MYCIN-style certainty factors: rules carry confidences,
conjunction is the G\"odel t-norm, and co-supporting rules combine through
a named aggregation operator (noisy-OR by default, provably MYCIN's
parallel-combination rule). Three properties distinguish it from its
lineage and from modern statistical-relational systems: rule confidences
are \emph{context-conditioned}, with a log-odds mode whose additive
evidence weights never saturate; rule and context weights are
\emph{learnable} --- on binary facts the engine's log-odds path reduces
exactly to logistic regression over rule firings, so training is convex
maximum likelihood; and explanations are \emph{traces whose fidelity is
measured} with deletion metrics against a reversed-attribution control.
The calculus carries a proved convergence proposition corroborated by
property-based tests, the engine is
a few hundred lines of dependency-free Python, and every reported number
regenerates byte-for-byte from committed scripts. Real-data studies on the
ULB credit-card and Bao et al.\ accounting-fraud datasets are included:
near-frontier ranking and a measured calibration/early-precision trade
on the former, and
pre-registered context ablations on the latter whose verdict --- no
reliable benefit --- we report as such.
\end{abstract}

\section{Introduction}
In audit, compliance, and clinical settings, a model that cannot show a
person why it concluded something is often unusable regardless of
accuracy; the case for inherently interpretable models over post-hoc
explanations is well established~\cite{rudin2019}. Systems with strong
probabilistic semantics exist (Section~\ref{sec:related}), but their
inference is opaque to practitioners and their rule weights are fixed at
inference time. PLA explores the opposite corner of the design space: a
calculus simple enough to trace by hand, honest about not being a
probability distribution, whose weights adapt to context and are learned
from data.

\textbf{Contributions.} (1)~An explicit confidence calculus whose formal
definitions execute as tests and whose fixpoint convergence is proved
for the monotone aggregators and corroborated by property tests on
1{,}050 random cyclic rule systems. (2)~Context-conditioned,
learnable rule weights: additive log-odds context evidence with convex
maximum-likelihood fitting that recovers planted ground truth. (3)~Measured
explanation fidelity: trace attributions out-score a reversed control on
deletion metrics. (4)~A reproducibility discipline: a zero-dependency
engine whose every reported number regenerates byte-for-byte, with each
claim mapped to a verifying artifact in the repository.

\section{Related work}\label{sec:related}
ProbLog attaches probabilities to logic programs under the distribution
semantics, computing exact possible-world marginals via knowledge
compilation and supporting conditioning on evidence --- an operation PLA
deliberately lacks~\cite{deraedt2007,fierens2015}; aProbLog generalizes
the semantics to commutative semirings~\cite{kimmig2011} --- a natural
reference point, though PLA's own pair is not a semiring instance
(G\"odel conjunction fails to distribute over noisy-OR); DeepProbLog and
Scallop add neural predicates and differentiable provenance
reasoning~\cite{manhaeve2018,li2023}. In weighted-logic systems --- Markov
Logic Networks~\cite{richardson2006}, Probabilistic Soft
Logic~\cite{bach2017}, Logical Neural Networks~\cite{riegel2020} --- rule
weights are static at inference time: adapting rule reliability to a
changing operational context is retraining's job, not the semantics'.
That gap, context-conditioned weights inside an interpretable calculus, is
the one PLA targets, connected to concept-drift
adaptation~\cite{gama2014}. All of these systems ship as Python packages;
PLA claims no language-based novelty. Its lineage is instead MYCIN's
certainty factors~\cite{shortliffe1975} with the G\"odel
t-norm~\cite{zadeh1965}; Heckerman showed the CF calculus is
probabilistically coherent only under restrictive
assumptions~\cite{heckerman1986,heckerman1992}. PLA's response: drop the
probability claim, name the independence stance as a per-KB operator
choice, and provide log-odds modes matching the coherent likelihood-ratio
special case.

\section{The calculus}\label{sec:calculus}
Propositions carry confidences $c(x)\in[0,1]$; base facts have $c=1$. A
rule $(A_1\wedge\dots\wedge A_n \rightarrow B,\; p,\; w)$ contributes
\begin{equation}
\text{candidate} \;=\; p_{\text{adj}} \cdot \min_i c(A_i),
\end{equation}
and candidates for the same head fold through a per-KB aggregator:
$\max$, noisy-OR $1-(1-e)(1-n)$ (identically MYCIN's parallel
combination), capped sum, or logit pooling
$\sigma(\operatorname{logit} e + \operatorname{logit} n)$. Context
adjustment has two modes: legacy multiply-and-cap, kept for
compatibility (the cap destroys information: distinct strong-evidence
combinations saturate identically at $1.0$), and the log-odds mode
\begin{equation}
p_{\text{adj}} \;=\; \sigma\bigl(\operatorname{logit}(p) +
\textstyle\sum_{v\,\text{active}} w(v)\bigr),
\end{equation}
in which strengthening never saturates and weakening is symmetric.
For the three monotone aggregators ($\max$, noisy-OR, capped sum),
inference iterates a monotone round operator on the lattice $[0,1]^V$;
the trajectory from the base facts is pointwise non-decreasing and
bounded, and the operators are continuous, so the limit is the Kleene
least fixpoint. (Logit pooling
is deliberately non-monotone --- sub-$\tfrac12$ candidates are negative
evidence --- and is used on the acyclic learning path.) These properties
are corroborated by property tests on 1{,}050 random cyclic systems, with a closed-form
regression test for slow high-probability cycles. The symbolic layer
decides definite-clause entailment by linear forward chaining
(differential-tested against a truth-table oracle; a 200-symbol KB
answers in well under a second) and gates confidences in hard/soft modes.
PLA has \emph{no} distribution semantics: no possible-world measure, no
conditioning, and aggregation chooses an independence stance --- users
needing those guarantees should use ProbLog.

\section{Learning}\label{sec:learning}
On flat rule sets over crisp binary facts, the log-odds context mode composed with logit pooling
collapses to $\hat{y} = \sigma\bigl(\sum_{r\,\text{fired}} z_r\bigr)$ with
$z_r = \operatorname{logit}(p_r) + \sum w_r(v)$ over active context
variables: logistic regression whose features are rule firings. The
binary cross-entropy is therefore convex in the parameters; full-batch
gradient descent decreases it monotonically (asserted over 200 epochs in
the test suite), and fitted parameters export back into the engine with
predictions matching to $10^{-9}$. On data sampled from a known model in
this family ($n{=}3{,}000$, three seeds), the learner recovers rule
weights within $0.09$ and the context weight within $0.29$.

\section{Tool and reproducibility}\label{sec:tool}
PLA ships as a pip-installable package with no core dependencies: a CLI
(\texttt{pla scenario.json [context-set]}), a Flask REST API behind an
optional extra, a fail-loud JSON scenario format (malformed context
shapes, unknown context sets, and undeclared variables raise instead of
silently applying no adjustment), and per-rule explanation traces as the
primary output. The repository enforces an unusual discipline for a
research prototype: 93 tests run in CI on Python 3.9--3.12 (one
self-skips off-network), every formula
in the semantics document is extracted and executed by the suite, the
head-to-head and fidelity result files are regenerated byte-for-byte by
their scripts, and the paper's claims each map to a verifying artifact.
Claims about competitor capabilities (ProbLog's marginals and evidence
conditioning; Bayesian-network posteriors) are themselves verified by
tests against the installed packages.

\section{Development results}\label{sec:results}
Numbers below are development results on a seeded synthetic sample with
planted logistic ground truth ($n{=}2{,}000$) that validate the harness;
the repository also carries two real-data studies (regenerated
byte-for-byte the same way), summarized at the end of this section. Table~\ref{tab:headtohead} compares PLA with static
empirical-precision rules, PLA with learned weights, and four baselines.
Reported as measured: with the standard logistic intercept included,
the learned variant matches static ranking (AUC $0.7227$ vs.\ $0.7217$)
and wins calibration decisively (log-loss $0.3420$ vs.\ $1.0908$) ---
raw empirical precisions are badly overconfident when several rules fire
together, and maximum likelihood repairs exactly that. Both PLA
variants trail the raw-feature baselines, as expected when the true model
is logistic in the raw features. The ProbLog row executes the same rules
under distribution semantics on the same examples and agrees with
PLA's noisy-OR fold to $7.1\times 10^{-7}$ at worst.

\begin{table}[t]
\centering
\caption{Head-to-head on the synthetic development sample (30\% test).
Generated by \texttt{scripts/run\_experiments.py}; reproduced
byte-for-byte from scratch.}
\label{tab:headtohead}
\begin{tabular}{lrrrr}
\toprule
Model & AUC & Brier & log-loss & ECE\\
\midrule
logistic regression (raw features) & 0.9746 & 0.0413 & 0.1390 & 0.0298\\
gradient boosting (raw features) & 0.9171 & 0.0883 & 0.3750 & 0.0808\\
ProbLog on PLA rules (subsample) & 0.7416 & 0.1327 & 1.4447 & 0.0704\\
pgmpy naive Bayes (propositions) & 0.7228 & 0.0993 & 0.3405 & 0.0288\\
PLA static (empirical precisions) & 0.7217 & 0.1138 & 1.0908 & 0.0952\\
PLA learned (logit path + bias) & 0.7227 & 0.0999 & 0.3420 & 0.0290\\
\bottomrule
\end{tabular}
\end{table}

Trace fidelity is measured with deletion metrics in the faithfulness
framing of Jacovi and Goldberg~\cite{jacovi2020}: comprehensiveness (drop
in prediction when the trace's top rule's facts are removed) and
sufficiency (gap when only they are kept), against a reversed-attribution
control. For both variants the faithful trace out-scores the control on
comprehensiveness ($0.2222$ vs.\ $0.1864$ static; $0.0939$ vs.\ $0.0807$
learned): the traces name factors that actually drive predictions --- a
measured property, not an asserted one (rule-level variants, log-odds
scales, and a random control are in the repository).

\paragraph{Real data.} The same scripts run on two real datasets. On the
ULB credit-card data (284{,}807 transactions, 492 frauds) an 11-rule PLA
model reaches AUC $0.9687$ (static) and $0.9625$ (learned, ECE $0.0022$)
against $0.9796$ for raw-feature logistic regression; average precision
exposes the trade behind that near-tie --- static $0.7258$ vs.\ learned
$0.5361$ --- so learning buys calibration at a measured early-ranking
cost. And the
interceptless learner's \emph{inverted} AUC of $0.076$, repaired to
$0.9625$ by the standard intercept, is a model-class lesson the logistic
reduction made obvious. On the Bao et al.\ accounting data~\cite{bao2020},
tail rules mined before the temporal boundary barely beat chance after it
(AUC $\approx 0.53$; raw-ratio logistic regression $0.68$--$0.70$), and
pre-registered context ablations --- including a design trained across
the 2003 SOX regime change so the context weight is identifiable ---
find no practically useful improvement from context conditioning on
the pre-specified AUC endpoint (paired $\Delta$AUC $0.0011$
[$-0.0071$, $0.0140$] under the SOX design; a post-hoc
average-precision gain on credit card awaits confirmation): an honest
negative, reported per its kill criterion.

\section{Limitations and roadmap}
PLA has no joint measure and no conditioning; its language is
propositional and negation-free (constraint-mode gating is a documented
no-op until negation lands); rule-shaped models pay a visible
discretization tax on continuous signal; the accounting-fraud study
simplifies the original protocol (complete-case deletion, an
approximation of its serial-fraud handling, a coarse rule vocabulary),
so its negative result binds that protocol, not the hypothesis space; and the usability claims
await a human-subject study. The roadmap: richer rule mining and the
full protocol of~\cite{bao2020} for the drift study, log-odds-scale
fidelity metrics for calibrated rare-event models, then the
applied-journal path.

\begin{thebibliography}{19}
\small
\bibitem{deraedt2007} L.~De Raedt, A.~Kimmig, H.~Toivonen.
ProbLog: a probabilistic Prolog and its application in link discovery.
\emph{IJCAI}, 2007.
\bibitem{fierens2015} D.~Fierens et al.
Inference and learning in probabilistic logic programs using weighted
Boolean formulas. \emph{TPLP} 15(3), 2015.
\bibitem{kimmig2011} A.~Kimmig, G.~Van den Broeck, L.~De Raedt.
An algebraic Prolog for reasoning about possible worlds. \emph{AAAI}, 2011.
\bibitem{manhaeve2018} R.~Manhaeve et al.
DeepProbLog: neural probabilistic logic programming. \emph{NeurIPS}, 2018.
\bibitem{li2023} Z.~Li, J.~Huang, M.~Naik.
Scallop: a language for neurosymbolic programming. \emph{PLDI}, 2023.
\bibitem{richardson2006} M.~Richardson, P.~Domingos.
Markov logic networks. \emph{Machine Learning} 62, 2006.
\bibitem{bach2017} S.~H.~Bach, M.~Broecheler, B.~Huang, L.~Getoor.
Hinge-loss Markov random fields and probabilistic soft logic.
\emph{JMLR} 18(109), 2017.
\bibitem{riegel2020} R.~Riegel et al. Logical neural networks.
arXiv:2006.13155, 2020.
\bibitem{gama2014} J.~Gama et al. A survey on concept drift adaptation.
\emph{ACM Computing Surveys} 46(4), 2014.
\bibitem{shortliffe1975} E.~H.~Shortliffe, B.~G.~Buchanan.
A model of inexact reasoning in medicine.
\emph{Mathematical Biosciences} 23, 1975.
\bibitem{zadeh1965} L.~A.~Zadeh. Fuzzy sets.
\emph{Information and Control} 8(3), 1965.
\bibitem{heckerman1986} D.~Heckerman.
Probabilistic interpretations for MYCIN's certainty factors.
\emph{Uncertainty in Artificial Intelligence 1}, North-Holland, 1986.
\bibitem{heckerman1992} D.~Heckerman, E.~H.~Shortliffe.
From certainty factors to belief networks.
\emph{Artificial Intelligence in Medicine} 4(1), 1992.
\bibitem{rudin2019} C.~Rudin.
Stop explaining black box machine learning models for high stakes
decisions and use interpretable models instead.
\emph{Nature Machine Intelligence} 1, 2019.
\bibitem{jacovi2020} A.~Jacovi, Y.~Goldberg.
Towards faithfully interpretable NLP systems: how should we define and
evaluate faithfulness? \emph{ACL}, 2020.
\bibitem{bao2020} Y.~Bao, B.~Ke, B.~Li, Y.~J.~Yu, J.~Zhang.
Detecting accounting fraud in publicly traded U.S. firms using a
machine learning approach. \emph{Journal of Accounting Research} 58(1),
2020.
\end{thebibliography}

\end{document}
